%% ---------------------------------------------------------------------
%% chem-dcat-ap-querying.tex
%% What a Query Result Means: Plan-Level Answering over Federated
%% Catalysis Data
%% ---------------------------------------------------------------------
\documentclass[11pt,a4paper]{article}
\usepackage[utf8]{inputenc}
\usepackage[T1]{fontenc}
\usepackage{amsmath,amssymb,amsthm,mathtools}
\usepackage{graphicx}
\usepackage{booktabs}
\usepackage{array}
\usepackage{enumitem}
\usepackage{float}
\usepackage{listings}
\usepackage{xcolor}
\usepackage{siunitx}
\DeclareSIUnit{\Molar}{M}
\usepackage{microtype}
\usepackage{lmodern}
\usepackage[margin=2.4cm]{geometry}
\usepackage[hidelinks]{hyperref}
\usepackage{cleveref}
\usepackage{natbib}
\crefname{lstlisting}{Listing}{Listings}
\Crefname{lstlisting}{Listing}{Listings}
\usepackage{stmaryrd}

\theoremstyle{plain}
\newtheorem{theorem}{Theorem}[section]
\newtheorem{lemma}[theorem]{Lemma}
\newtheorem{proposition}[theorem]{Proposition}
\newtheorem{corollary}[theorem]{Corollary}
\theoremstyle{definition}
\newtheorem{definition}[theorem]{Definition}
\newtheorem{assumption}[theorem]{Assumption}
\newtheorem{construction}[theorem]{Construction}
\newtheorem{principle}[theorem]{Principle}
\theoremstyle{remark}
\newtheorem{remark}[theorem]{Remark}
\newtheorem{example}[theorem]{Example}
\newtheorem{observation}[theorem]{Observation}

% ---- notation, shared with the federated-query technical report -------
\newcommand{\Plan}{\mathcal{P}}
\newcommand{\Step}{\sigma}
\newcommand{\Src}{\mathsf{S}}
\newcommand{\Capset}{\mathrm{cap}}
\newcommand{\Req}{\mathrm{req}}
\newcommand{\Low}{\mathrm{low}}
\newcommand{\Ext}{\mathrm{ext}}
\newcommand{\Res}{\mathcal{R}}
\newcommand{\Bind}{\Gamma}
\newcommand{\Verd}{\mathrm{v}}
\newcommand{\vans}{\textsf{answered}}
\newcommand{\vempty}{\textsf{empty}}
\newcommand{\vsurf}{\textsf{unsupported}}
\newcommand{\vtime}{\textsf{exhausted}}
\newcommand{\vref}{\textsf{refused}}
\newcommand{\vstarve}{\textsf{starved}}
\newcommand{\bmdl}{\textsf{blocked-by-model}}
\newcommand{\bcorp}{\textsf{blocked-by-corpus}}
\newcommand{\beng}{\textsf{blocked-by-engine}}
\newcommand{\bbud}{\textsf{blocked-by-budget}}
\newcommand{\Univ}{\mathcal{U}}
\newcommand{\Nspace}{\mathcal{N}}
\newcommand{\idm}{\mathrm{id}}
\newcommand{\Rset}{\mathbb{R}}
\newcommand{\Nset}{\mathbb{N}}
\newcommand{\Bud}{B}
\newcommand{\cost}{c}
\newcommand{\HFQ}{\textsc{hfq}}
\newcommand{\Feat}{\mathcal{F}}
\newcommand{\Zset}{\mathbb{Z}}

\lstdefinestyle{hfq}{
  basicstyle=\ttfamily\footnotesize,
  keywordstyle=\bfseries,
  commentstyle=\itshape\color{gray!70!black},
  stringstyle=\color{black},
  showstringspaces=false,
  breaklines=true,
  frame=single,
  framesep=4pt,
  rulecolor=\color{gray!45},
  xleftmargin=6pt,
  morekeywords={let,from,ask,with,in,within,emit,budget,requests,union,
    intersect,join,on,filter,where,map,via,ladder,over,provenance,as,
    extension,of,because,divergence}
}
\lstdefinestyle{sparql}{
  basicstyle=\ttfamily\footnotesize,
  keywordstyle=\bfseries,
  commentstyle=\itshape\color{gray!70!black},
  showstringspaces=false,
  breaklines=true,
  frame=single,
  framesep=4pt,
  rulecolor=\color{gray!45},
  xleftmargin=6pt,
  morekeywords={SELECT,DISTINCT,WHERE,VALUES,OPTIONAL,FILTER,SERVICE,
    NOT,EXISTS,MINUS,PREFIX,BIND,GROUP,BY,ORDER,LIMIT,UNION}
}

\title{What a Query Result Means:\\
Plan-Level Answering over Federated Catalysis Data}
\author{Kundai Farai Sachikonye}
\date{September 2026}

\begin{document}
\maketitle

\begin{abstract}
The standard route from a chemical question to an answer is: build an
ontology, materialise a knowledge graph, apply a reasoner, and issue a
SPARQL query. This paper reports what happens when thirteen questions
drawn from working biocatalysis practice are put instead to a
\emph{plan} language with an explicit capability calculus, and argues
that the difference is not one of convenience but of what the returned
document is able to say.

The central observation is negative and precise. A SPARQL result set
carries one bit of failure information: the solution sequence is either
non-empty or empty. Four operationally distinct situations collapse onto
that empty result --- the corpus genuinely contains no such thing; a
source cannot evaluate the constraint at all; an upstream step yielded
nothing so a downstream step was never reachable; the request was
declined. We show that no rewriting of the query recovers the
distinction, because the interface has no channel to carry it. We then
exhibit a six-valued verdict assigned by ordered rules over each plan
step, a capability check that refuses ill-posed plans \emph{before any
request is issued}, and a fourth answer shape --- an admissibility
statement --- for the case where every step succeeds and the plan
nonetheless answers a different question than the one asked.

All thirteen questions execute. Five are Doerr's biocatalysis
questions; eight are the generic Chem-DCAT-AP dataset-discovery
patterns. Eight return rows. Two are refused statically at
\texttt{requests\_issued} $= 0$, one of them for data that demonstrably
exists in the federation. Two return rows under an explicit
admissibility caveat naming the gap. One returns an \vempty{} that is
reported as \vempty{} and a downstream \vstarve{} that is reported as
\vstarve{}, which is the whole point. We give the counterfactual
falsification of every constraint --- including one that caught a real
defect in our own plan, where a checked constraint was never on the path
from the question to the payload and the plan returned the right rows
for the wrong reason. Sixteen validation checks over the compiler and
executor hold. We are explicit about what they do not establish.
\end{abstract}

\tableofcontents

% =====================================================================
\section{The question this paper is about}
\label{sec:intro}

A biocatalysis group keeps its experimental records in a structured
form. A colleague asks:

\begin{quote}
\emph{Which biocatalyst, originating from a bacterium and not a
eukaryote, catalyses the transamination of benzylethylamine and does not
have a cysteine in its protein sequence?}
\end{quote}

This is a reasonable question and it has, in the group's data, a
determinate answer. The received method for obtaining it is well
established and works: express the domain in an ontology, materialise
the records as RDF, run a reasoner to close the model under the
axioms, and issue a SPARQL query against the result
\citep{Baader2003,Grau2008,Harris2013}. The Chem-DCAT-AP effort applies
exactly this recipe to catalysis data, extending DCAT-AP
\citep{DCATAP2015,DCAT2020} with the chemistry-specific vocabulary the
domain needs, and it does so competently.

We are not going to argue that this does not work. It does. We are
going to argue that it works in a way that cannot tell you when it has
not worked, and that for laboratory data --- which is sparse,
heterogeneous, and unevenly recorded by construction --- that
limitation is not incidental.

\subsection{A single observation}

Consider what the SPARQL engine hands back when the question above
returns nothing. It hands back an empty solution sequence. Now consider
the situations that produce an empty solution sequence here:

\begin{enumerate}[label=(\alph*),leftmargin=2.2em]
\item Every transaminase in the corpus that acts on benzylethylamine has
  a cysteine. The question has a genuine negative answer.
\item The sequence store is a REST service that returns a FASTA record
  by accession and cannot be asked \emph{which of your records lack
  residue X}. The constraint was never evaluated.
\item No transaminase in the corpus is recorded as bacterial, because
  organism attribution lives in a source that was not consulted. The
  cysteine filter had nothing to filter.
\item The endpoint declined the request --- rate limit, timeout,
  malformed federation clause.
\end{enumerate}

These are four different facts about the world and about the query
apparatus, and they call for four different next actions: publish the
negative result; find a different sequence source; add the organism
source to the federation; retry. The interface returns the same thing in
all four cases.

\begin{observation}[the one-bit interface]
\label{obs:onebit}
The SPARQL result-set interface has exactly two failure-relevant states:
non-empty and empty. Situations (a)--(d) are distinguishable facts about
an execution that are mapped onto one state. No rewriting of the query
text changes this, because the distinction is not a property of the
query --- it is a property of the execution, and the interface has no
field in which to put it.
\end{observation}

This is not a criticism of any implementation. It is a statement about
what the abstraction was designed to convey. SPARQL was designed to
express \emph{what set of solutions do the data entail}; it answers that
question well. It was not designed to convey \emph{how did the attempt
to obtain them go}, and asking it to is a category error.

The consequence, however, is borne by the chemist. In a domain where the
data are complete --- a well-curated single graph, closed-world by
construction --- case (a) is overwhelmingly the likely reading of an
empty result and the ambiguity is harmless. Laboratory catalysis data is
the opposite case. Records are partial by nature: a run that did not
record its buffer is not a run that used no buffer. Sources are
heterogeneous: sequences live in UniProt \citep{UniProt2023}, kinetics
in SABIO-RK \citep{Wittig2018}, enzyme classification in BRENDA
\citep{Chang2021}, compound identity in ChEBI \citep{Hastings2016}, and
the group's own provenance in its own store. Under those conditions
(b), (c), and (d) are not edge cases, they are the common case, and an
empty result that cannot be read is the normal outcome.

\subsection{What we propose instead}

We keep the ontology work. Nothing here argues against having a
vocabulary; you cannot ask about a \emph{catalyst} without a term for
one. What we change is the layer above: instead of compiling the
question into a single query and reading the result set, we compile it
into a \emph{plan} --- an ordered sequence of steps, each bound to a
named source, each carrying a declared set of the capabilities it
requires --- and we execute the plan against a federation whose sources
declare what they can do.

Three things follow, and they are the substance of this paper.

\begin{enumerate}[leftmargin=2.2em]
\item \textbf{The result carries a verdict per step, drawn from six
  values, not two} (\Cref{sec:verdicts}). Situations (a)--(d) above
  receive distinct verdicts, assigned by ordered rules, with a
  \emph{blocker} naming what stood in the way.

\item \textbf{Plans that cannot be answered are refused before contact}
  (\Cref{sec:refusal}). If a step requires a capability its source does
  not declare, the plan is rejected statically and no request is issued
  --- including the requests for the parts that would have succeeded. We
  demonstrate this at \texttt{requests\_issued} $= 0$ on two questions,
  one of which asks for data that provably exists in the federation.

\item \textbf{Some questions receive an answer and an admissibility
  statement} (\Cref{sec:admissibility}). A plan can return the correct
  rows at every step and still be answering a different question,
  because the question's logical form and the corpus's do not match.
  Nothing in a verdict can say so, so we say it separately, and we make
  the plan author name which kind of gap it is.
\end{enumerate}

\subsection{For the chemist, and for the computer scientist}

This paper has two audiences and they want different things from it. We
have tried to serve both rather than splitting the difference.

To the chemist: the claim you should test is in \Cref{sec:questions},
where five real questions are worked end to end. The claim is not that
we answer more of them --- we answer eight of thirteen with rows, which
a well-built SPARQL stack over the same data would also do. The claim is
that for the five we do not answer, the document we return tells you
\emph{why not}, in terms that distinguish a fact about your chemistry
from a fact about your infrastructure. \Cref{sec:q4} is the one to read
first: it is a question about expected products, and the answer is
that the corpus contains the reaction and the conditions and still
cannot support the inference, for a reason that a reasoner would
\emph{hide} rather than expose.

To the computer scientist: the technical content is the capability
calculus of \Cref{sec:calculus}, the verdict rules R1--R6 of
\Cref{sec:verdicts}, and the honest-declaration assumption of
\Cref{sec:honesty} --- which is the largest gap in the design and which
we do not close. The relationship to prior federation work
\citep{Wiederhold1992,Levy1996,Haas1997,Schwarte2011,Saleem2016} is
that source-capability description is old and well studied; what is new
here is treating the \emph{failure} of a capability match as a
first-class output rather than as a planner-internal fact, and the
observation that this is what the empty-result ambiguity actually needs.
The connection to the nulls and certain-answers literature
\citep{Codd1979,Imielinski1984,Libkin2016} is close and we make it
explicit in \Cref{sec:related}.

\subsection{What this paper does not claim}

We state the limits at the front rather than burying them.

\begin{enumerate}[label=(\roman*),leftmargin=2.2em]
\item The capability declarations are \emph{data written by the adapter
  author}. Nothing verifies them. An over-declaring source produces
  exactly the silent wrong answer this paper exists to prevent
  (\Cref{sec:honesty}).
\item The validation suite tests the compiler and the executor against
  the theory. It does not test the theory against the world
  (\Cref{sec:validation}).
\item The evaluation is over a local fixture built to mirror the
  structure of the group's data, not over the live endpoint. This is
  deliberate and we explain the ethics of it in \Cref{sec:noprobing}.
\item Eight of thirteen questions return rows. We are not claiming
  higher recall than a conventional stack; on this corpus we would
  expect the same recall. The claim is about the other five.
\end{enumerate}

% =====================================================================
\section{Plans, sources, capabilities}
\label{sec:calculus}

This section fixes the objects. A reader who wants the chemistry first
may skip to \Cref{sec:questions} and return here.

\subsection{The move: a question is a plan, not a query}

The received pipeline compiles a question into one query and submits it.
The query is a single syntactic object whose evaluation is the
endpoint's business, and the planner's decisions --- which source, in
what order, with what budget --- are either absent (one graph) or
internal to a federation engine that does not report them.

We invert this. The user writes a \emph{plan}: an ordered list of named
steps, each of which names a source and a predicate, and may bind the
output of an earlier step. Queries are \emph{leaves}. The user never
writes one. The lowering from a step to a concrete request in the target
language is the adapter's job, and the lowered form is retained in the
output so that what was actually sent can be read.

\begin{lstlisting}[style=hfq,caption={A plan. Every step names its
source; no SPARQL appears anywhere.},label={lst:plan}]
plan mark_q1 {
  budget 12 requests

  let enzymes = from RXN
      ask catalysts_of("RXN:TA-benzylethylamine")
      within 20

  let bacterial = from TAX
      ask typed_as("TAX:bacteria")
      with ?e in enzymes
      within 20

  let no_cys = from SEQ
      ask lacks_residue("C")
      with ?e in candidates
      within 20

  emit no_cys with provenance
}
\end{lstlisting}

The reason for the inversion is not ergonomics. It is that the
information we want in the output --- which source could not do what ---
exists only at the granularity of a step. A single query has one
outcome; a plan has one per step, and it is at that granularity that
\Cref{obs:onebit} is escapable.

\subsection{Sources}

\begin{definition}[source]
\label{def:source}
A source is a quadruple $\Src = (\eta, \Capset, \Ext, \cost)$ where
$\eta$ is a namespace identifier, $\Capset \subseteq \Feat$ is the set of
capabilities the source declares, $\Ext$ is the extraction function
mapping a lowered request and its bound inputs to a result set, and
$\cost \in \Nset$ is the request cost charged against the budget.
\end{definition}

\begin{definition}[capability vocabulary]
\label{def:feat}
$\Feat$ is the fixed, finite set
\[
\Feat = \{\,\textsf{pattern},\ \textsf{path},\ \textsf{filter},\
\textsf{bind},\ \textsf{agg},\ \textsf{neg},\ \textsf{order},\
\textsf{regex},\ \textsf{lookup},\ \textsf{link},\ \textsf{batch}\,\}
\]
of eleven symbols. Each names an operation a source either can or cannot
perform: matching a triple pattern; traversing a path, including
backwards; applying a value restriction; accepting a set of bound values
as input; aggregating; evaluating negation; ordering; matching a regular
expression; retrieving a record by key; following a link; and accepting
a batch of keys in one request.
\end{definition}

The vocabulary is deliberately small and deliberately closed. Its
purpose is not to describe a source completely --- it cannot --- but to
be coarse enough that an adapter author can answer each of eleven
yes/no questions honestly. A finer vocabulary would be more expressive
and less reliably answered, and the failure mode of a wrong answer here
is severe (\Cref{sec:honesty}).

\begin{definition}[requirement]
\label{def:req}
Each predicate $p$ carries a requirement $\Req(p) \subseteq \Feat$. The
requirement of a step $\Step$ is $\Req(\Step) = \Req(p) \cup
\{\textsf{bind}\}$ if $\Step$ binds any earlier step, and $\Req(p)$
otherwise.
\end{definition}

\begin{definition}[well-capability]
\label{def:wellcap}
A plan $\Plan$ over sources $\{\Src_i\}$ is \emph{well-capable} iff for
every step $\Step$ with source $\Src$,
$\Req(\Step) \subseteq \Capset(\Src)$.
\end{definition}

\subsection{Honest declaration, and its asymmetry}
\label{sec:honesty}

The critical property of \Cref{def:source} is what happens when the
declaration is wrong, and it is not symmetric.

\begin{proposition}[under-declaration is safe, over-declaration is not]
\label{prop:honest}
Let $\Src$ have true capability set $C^{\ast}$ and declared set
$\Capset$.
\begin{enumerate}[label=(\alph*),leftmargin=2.2em]
\item If $\Capset \subsetneq C^{\ast}$, then some plans that the source
  could in fact serve are refused. The refusal is visible, names the
  missing capability, and is repairable by amending the declaration. No
  wrong answer is produced.
\item If $\Capset \supsetneq C^{\ast}$, then a plan passes the static
  check, the request is issued, the source returns something --- a
  well-formed result set --- and that result set is wrong in a way no
  field of the output records. The failure is invisible.
\end{enumerate}
\end{proposition}

\begin{proof}
(a) The check of \Cref{def:wellcap} is a subset test computed before
execution; a missing element yields a refusal naming
$\Req(\Step) \setminus \Capset(\Src)$. (b) The check passes by
construction, so no refusal is emitted; the adapter's lowering produces
a request the source accepts under a reading the adapter did not intend,
and the extraction returns whatever that reading yields. Nothing in the
pipeline compares the returned rows against the intended semantics,
because there is no independent statement of the intended semantics to
compare against.
\end{proof}

\begin{remark}[the honesty assumption, and the largest gap]
\label{rem:honesty}
\Cref{prop:honest}(b) is the single largest gap in this design and we do
not close it. The capability declaration is data written by a human. A
declaration claiming \textsf{neg} for a source that silently ignores
negation produces precisely the confident wrong answer that motivates
the whole apparatus. The design's answer is architectural rather than
technical: the declaration is small (eleven bits), local, reviewable,
and versioned alongside the adapter, so that it is the kind of artefact
a reviewer can actually check. That is a mitigation, not a proof, and a
reader should weigh it as such.
\end{remark}

A concrete instance is worth stating because it inverts the usual
intuition.

\begin{proposition}[a supported operation that must be withheld]
\label{prop:under}
Let a source expose a SPARQL endpoint with a server-side
materialisation cap: aggregate queries over result sets exceeding the
cap return a well-formed aggregate computed over a truncated set,
with an HTTP 200 and no warning. Then an honest declaration must exclude
\textsf{agg} from $\Capset$, even though SPARQL supports aggregation and
the request succeeds.
\end{proposition}

\begin{proof}
The declaration in \Cref{def:source} asserts that the source can
\emph{perform} the operation, not that it will \emph{accept} a request
naming it. A truncated aggregate is not the aggregate; it is a
different number returned under the same name. Including \textsf{agg}
places the deployment in case (b) of \Cref{prop:honest}, whose failure
is invisible by construction. Excluding it places the deployment in case
(a), whose failure is visible and repairable --- the plan is refused,
and the plan author computes the aggregate client-side over rows the
source \emph{can} return faithfully.
\end{proof}

This is the sharpest single point we can make to a chemist about why
the capability vocabulary is not bureaucratic overhead. Under the
conventional pipeline the truncated aggregate is a number that appears
in a table in a paper. Under this one, the plan does not run.

\subsection{Budget and allocation}

Steps cost requests, and a plan declares a budget. When the budget binds,
something must give, and \emph{how} it gives is a semantic decision, not
an engineering one.

\begin{definition}[yield]
\label{def:yield}
The yield of step $i$ is a function $\gamma_i : \Rset_{\geq 0} \to
\Rset_{\geq 0}$ from effort, measured in requests, to expected returned
cardinality. We assume $\gamma_i$ is non-decreasing, continuously
differentiable, strictly concave, with $\gamma_i(0) = 0$.
\end{definition}

\begin{theorem}[water-filling allocation]
\label{thm:allocation}
Under \Cref{def:yield}, the allocation problem
$\max \sum_i \gamma_i(e_i)$ subject to $\sum_i \cost_i e_i \leq \Bud$,
$e_i \geq 0$, has a unique optimum $e^{\ast}$ characterised by a single
scalar shadow price $p^{\ast} \geq 0$ with
$\gamma_i'(e_i^{\ast}) = p^{\ast} \cost_i$ for every $i$ with
$e_i^{\ast} > 0$, and $\gamma_i'(0) \leq p^{\ast} \cost_i$ otherwise.
\end{theorem}

\begin{corollary}[no dropping rule]
\label{cor:nodrop}
No step is excluded from the allocation by heuristic, priority order, or
blacklist. A step receives zero effort only when its marginal yield at
zero is below the shadow price, which is a statement about the step and
the budget jointly, and is reported as such.
\end{corollary}

\Cref{cor:nodrop} matters for the same reason the rest of this does. A
federation engine that drops a subquery under load returns a result set
that is missing rows for a reason the caller cannot see. Here, a step
that received no effort is reported with a \bbud{} blocker, and the
caller knows the difference between \emph{there were none} and
\emph{we did not look}.

\begin{remark}[where concavity fails]
\label{rem:concavity}
A \textsf{lookup} against a keyed REST service is all-or-nothing: one
request returns the record or nothing, so $\gamma$ is a step function
and not concave. \Cref{thm:allocation} does not apply to it. Such steps
are charged first at their fixed cost and the remainder is optimised
over the concave steps. We note this rather than pretending the
assumption is universal.
\end{remark}

% =====================================================================
\section{Six verdicts, and why two are not enough}
\label{sec:verdicts}

\subsection{The rules}

Each executed step receives exactly one verdict, assigned by the first
matching rule of an ordered list. Ordering is what makes the assignment
total and deterministic.

\begin{definition}[verdict rules]
\label{def:rules}
For a step $\Step$ with source $\Src$, bound inputs $\Bind$, allocation
$e$, and outcome $O$:
\begin{description}[leftmargin=3.2em,style=nextline]
\item[\textnormal{R1}] If $\Req(\Step) \not\subseteq \Capset(\Src)$:
  $\vsurf$, blocker \bmdl. The source cannot perform the operation.
\item[\textnormal{R2}] If $e = 0$ under a binding budget: $\vtime$,
  blocker \bbud. Effort was not allocated.
\item[\textnormal{R3}] If some bound input $\Bind_j$ is empty:
  $\vstarve$, blocker \bcorp. The step was unreachable; it was not
  attempted.
\item[\textnormal{R4}] If the source declined or errored: $\vref$,
  blocker \beng.
\item[\textnormal{R5}] If the request completed and returned no rows:
  $\vempty$, \emph{no blocker}.
\item[\textnormal{R6}] Otherwise: $\vans$, no blocker.
\end{description}
\end{definition}

\begin{proposition}[totality and determinism]
The rules assign exactly one verdict to every executed step.
\end{proposition}
\begin{proof}
R6 is unconditional, so the assignment is total; the rules are consulted
in order and the first match returns, so it is unique.
\end{proof}

\subsection{The absent blocker is the point}

R5 assigns no blocker, and this is a design decision that a reader
should push on rather than accept.

\begin{principle}[the blocker map is partial on purpose]
\label{prin:noblocker}
$\vempty$ carries no blocker, and the corresponding JSON field is
\emph{absent} rather than null. An empty result that completed is not
blocked by anything. The corpus was consulted, the constraint was
evaluated, and the extension is empty. That is a fact about the
chemistry.
\end{principle}

Every other verdict is a fact about the apparatus. Keeping the
distinction in the shape of the output --- a missing key, not a null
--- means a consumer that reads \texttt{blocker} and finds nothing has
been told something, and a consumer that reads \texttt{null} has been
told nothing and may believe it was told something.

This is the resolution of \Cref{obs:onebit}. Situations (a)--(d) of
\Cref{sec:intro} receive $\vempty$ (no blocker), $\vsurf$ (\bmdl),
$\vstarve$ (\bcorp), $\vref$ (\beng) respectively.

\begin{corollary}[the distinction is not recoverable downstream]
\label{cor:onebit}
No SPARQL query, however written, recovers the
$\vempty$ / $\vstarve$ / $\vsurf$ / $\vref$ distinction from a result
set, because the distinction is a property of the execution and the
result-set interface has no channel for it. Recovering it requires
changing the interface, which is what a plan-level verdict does.
\end{corollary}

\subsection{Starvation is the characteristic federated failure}

\begin{proposition}[starvation is reachable only under federation]
\label{prop:starve}
Over a single closed corpus, R3 cannot fire on a plan whose steps are
all satisfiable: an empty bound input is itself an empty extension, and
the composite is empty by R5. Under federation with a translation
between namespaces, a step may be unreachable because a \emph{mapping}
--- not the corpus --- yielded nothing for the identifiers at hand.
\end{proposition}

This is why $\vstarve$ has its own verdict rather than being folded into
$\vempty$. \emph{Nothing in the sequence store lacks a cysteine} and
\emph{no identifier we hold could be translated into the sequence
store's namespace} are different problems, and only the second is fixed
by adding a mapping.

We measure translation with two independent quantities.

\begin{definition}[retention and amplification]
\label{def:retamp}
For a translation map $\mu$ and a set $S$ of identifiers in its domain
namespace, the retention is $r_\mu(S) = |\{s \in S : \mu(s) \neq
\emptyset\}| / |S|$ and the amplification is $a_\mu(S) = |\bigcup_{s\in
S}\mu(s)| \big/ |\{s \in S : \mu(s) \neq \emptyset\}|$, undefined when
the denominator is zero. By convention $r_\mu(\emptyset) = 1$.
\end{definition}

\begin{proposition}
Retention and amplification are independent, and neither is determined
by output cardinality.
\end{proposition}
\begin{proof}
A map sending four of eight identifiers to one image each has
$r = 0.5$, $a = 1$, output $4$; a map sending two of eight to two images
each has $r = 0.25$, $a = 2$, output $4$. Equal outputs, different
$(r,a)$.
\end{proof}

Question Q1 (\Cref{sec:q1}) exercises this: its translation step from
enzyme identifiers to sequence accessions records $r = 0.75$, $a = 1.0$.
One in four candidate enzymes has no sequence record. A conventional
pipeline reports four candidates in and three out, and the missing one
is indistinguishable from a filtered one.

% =====================================================================
\section{Refusal before contact}
\label{sec:refusal}

\begin{corollary}[refuse before contact]
\label{cor:refuse}
If a plan is not well-capable, execution halts at the check stage and
$\texttt{requests\_issued} = 0$.
\end{corollary}

Two aspects of this are stronger than they may first appear, and both
are load-bearing.

\subsection{The whole plan, not the failing step}

\begin{principle}[ill-capability is a property of the plan]
\label{prin:whole}
When one step fails the capability check, the satisfiable steps do not
run either.
\end{principle}

The argument is not conservatism. Suppose a question imposes two
constraints, one evaluable and one not, and suppose the evaluable half
runs and emits its rows. The caller now holds a set of datasets that
satisfy constraint A and an \emph{unenforced} constraint B, and no field
in the payload distinguishes them from datasets satisfying both. The
partial result is more dangerous than no result, because it has the
shape of an answer. \Cref{sec:g8} exhibits this concretely.

\subsection{The refusal names the real obstacle}

\begin{principle}[name the obstacle]
\label{prin:refusal}
A refusal states which step, which source, which capabilities were
required, which were declared, and which were missing.
\end{principle}

\begin{lstlisting}[style=hfq,caption={The emitted refusal for G8.
\texttt{requests\_issued} is 0.},label={lst:refusal}]
{"refusal": {"plan": "dcat_g8",
  "outcome": "refused_statically",
  "reason": "ill-capability plan; no request was issued",
  "failures": [{"step": "uvvis", "source": "INST",
     "required": ["neg", "pattern"],
     "declared": ["batch", "bind", "link", "lookup"],
     "missing":  ["neg", "pattern"]}],
  "steps": []}}
\end{lstlisting}

A chemist reading this learns something actionable: the instrument
registry is a keyed record service, and the question requires asking it
to classify its own contents, which it cannot do. The fix is to add a
source that can, or to reformulate. Compare what the conventional
pipeline offers for the same situation: an empty result, and an
afternoon.

% =====================================================================
\section{When every step succeeds and the answer is still wrong}
\label{sec:admissibility}

This section contains the argument we consider most important, and it is
the one that is hardest to make to an audience that has internalised the
reasoner.

\subsection{The problem}

A plan can be $\vans$ at every step, return exactly the rows its
predicates specify, and answer a different question than the one asked.
The mismatch is not between a step and a source --- which is what the
verdict algebra measures --- but between the \emph{question's logical
form} and the \emph{corpus's}. Nothing in \Cref{def:rules} can detect
it, because every rule is about the fate of a request.

\begin{example}[substrate scope]
\label{ex:scope}
\emph{What is the substrate scope of BVMO-Y?} The corpus records two
substrates on which BVMO-Y was run. A plan retrieves both, every step
answers, and the emitted rows are correct. But \emph{scope} is a claim
about substrates the enzyme \emph{would} accept, including ones nobody
has tried. The returned set is the recorded extension; the question
asked for an intension. These differ, and no traversal of this corpus
closes the gap.
\end{example}

The gap in \Cref{ex:scope} is inductive: from observed instances to a
general claim. That step is chemistry --- structure-activity reasoning,
homology, mechanism --- and it is not inference over the recorded facts.

\subsection{Why a reasoner makes this worse}

Here is the sharpest divergence between the two methods, and it runs
against the intuition that more inference is more help.

\begin{proposition}[asserting the closure launders an assumption]
\label{prop:launder}
Let an ontology assert an axiom sufficient to derive substrate scope
from recorded substrates --- for instance, that the substrate class of
an enzyme is closed under a structural similarity relation. Then the
reasoner derives triples asserting acceptance of untested substrates,
and those derived triples are indistinguishable, in the materialised
graph and in any query result over it, from triples recording observed
acceptance.
\end{proposition}

\begin{proof}
Materialisation writes entailed triples into the same graph as asserted
ones, in the same vocabulary. A SPARQL query matching the substrate
predicate matches both. Distinguishing them requires provenance
annotation on every derived triple and a query discipline that consults
it, neither of which the closure operation provides by default.
\end{proof}

The situation is therefore: the honest answer is \emph{the corpus
records two substrates and cannot tell you the scope}; the reasoner's
answer is \emph{a list}, and the list is a mixture of observation and
assumption that looks homogeneous. The chemist who reads that list has
been given a prediction wearing the clothes of a record. This is not a
hypothetical failure mode; it is the ordinary consequence of
materialising a closure over sparse experimental data, and it is exactly
why \citet{Reiter1978,Reiter1980} treated closed-world assumption and
default inference as things to be marked rather than absorbed.

\subsection{The construct}

We add a fourth answer shape. It sits beside the verdict, not inside
it.

\begin{construction}[admissibility statement]
\label{con:admiss}
A plan may emit
\begin{center}
\texttt{emit X with provenance as extension of "TEXT" because GAP}
\end{center}
where $\textsf{GAP} \in \{\textsf{induction}, \textsf{vocabulary},
\textsf{conditions}\}$. The emitted document then carries an
\texttt{admissibility} record stating what was asked for, that what is
returned is the recorded extension, that
\texttt{answers\_question: false}, and a reason determined by the gap.
\end{construction}

Two features of \Cref{con:admiss} are deliberate.

First, the verdicts stay honest. The steps did answer; recording
\vref{} here would falsely suggest retrieval failed, and a reader could
not distinguish a corpus that lacks the rows from one that has them and
cannot generalise over them. So \texttt{answers\_question: false}
appears \emph{alongside} \texttt{verdict: answer}.

Second, the gap is named by the plan author from a closed set, not
generated. A free-text explanation is the same defect one level up: a
canned sentence that sounds like reasoning and is not checked. A closed
set of three is small enough to review and forces a claim about
\emph{which kind} of gap this is.

\begin{itemize}[leftmargin=1.6em]
\item \textsf{induction} --- the question asks about cases the corpus
  does not contain (\Cref{ex:scope}).
\item \textsf{vocabulary} --- the question uses a term the corpus does
  not record as such, though it may record something adjacent.
\item \textsf{conditions} --- the question asks for an outcome under
  conditions the corpus records for a \emph{different} configuration.
\end{itemize}

\Cref{sec:g6} shows that \textsf{vocabulary} covers two structurally
distinct cases --- a missing type on a subject the store holds, versus
an entirely absent relation --- which is a limitation of the current
three-way partition that we name rather than paper over.

% =====================================================================
\section{Five questions from a working group}
\label{sec:questions}

This section works the five biocatalysis questions end to end. They were
supplied by a practising group, not constructed by us, and we did not
alter them. Where a question cannot be answered we say so and say why;
where our first attempt was wrong we show the wrong version.

The federation comprises six sources, each declaring a capability set
under \Cref{def:source}: \textsf{RXN} (reactions and catalysis, a graph
source), \textsf{TAX} (organism attribution), \textsf{SEQ} (protein
sequences, a keyed record service), \textsf{ACT} (activities and runs,
a provenance source), \textsf{INST} (an instrument registry, keyed
records), and \textsf{DS} (dataset descriptions). One translation map,
$\mu : \textsf{enz} \to \textsf{seq}$, crosses the enzyme and sequence
namespaces.

\subsection{Q1: bacterial transaminase without cysteine}
\label{sec:q1}

\begin{quote}
\emph{Which biocatalyst, originating from a bacterium (not a
eukaryote), catalyses the transamination of benzylethylamine and does
not have a cysteine in its protein sequence?}
\end{quote}

This is one sentence and three sources. Catalysis is in \textsf{RXN};
organism attribution is in \textsf{TAX}; residue composition is in
\textsf{SEQ}, whose identifiers are accessions and not enzyme
identifiers. \Cref{lst:plan} gives the plan. Its execution:

\begin{table}[H]\centering
\begin{tabular}{llrrl}
\toprule
step & verdict & $n$ & $r_\mu$ / $a_\mu$ & note \\
\midrule
\texttt{enzymes}    & \vans & 5 & --- & catalysts of the transamination \\
\texttt{bacterial}  & \vans & 4 & --- & one eukaryotic entry excluded \\
\texttt{candidates} & \vans & 4 & --- & intersection \\
\texttt{keyed}      & \vans & 3 & $0.75$ / $1.0$ & namespace translation \\
\texttt{no\_cys}    & \vans & 2 & --- & residue filter \\
\bottomrule
\end{tabular}
\caption{Q1. Three requests issued, budget twelve, no early halt.}
\label{tab:q1}
\end{table}

Two enzymes are returned, and that is the answer. What we want to draw
attention to is the fourth row.

Four candidates entered the translation and three left it. The step
reports $r_\mu = 0.75$: one in four candidate enzymes has no sequence
record in the accession namespace. It also reports $a_\mu = 1.0$: no
enzyme mapped to more than one accession, so no candidate was silently
duplicated.

Under the conventional pipeline this step is a join, and a join that
loses a row loses it silently. The chemist sees four candidates in one
intermediate result and two in the final, and has no way to apportion
the attrition between \emph{filtered by the cysteine constraint} and
\emph{never reached the sequence store}. Those have different remedies.
The first is a finding. The second is a missing cross-reference, and it
is fixed by curation, not by chemistry.

\begin{remark}
This is the only step in all thirteen plans that carries map metrics,
because it is the only step that crosses a namespace. We state this
rather than leaving the impression that retention and amplification are
reported everywhere. Where there is no translation there is no
attrition to measure, and the fields are absent --- not zero.
\end{remark}

\subsection{Q2: buffer and pH of a specific methyl transfer}
\label{sec:q2}

\begin{quote}
\emph{Which buffer composition and pH was used in the biocatalytic
methyl transfer using methyl-transferase mt-X of Bacillus subtilis?}
\end{quote}

Five steps, five requests, all \vans: two \emph{B.\ subtilis} enzymes,
of which one is MT-X; one methyl-transfer reaction; two runs; two
condition records. The plan returns the buffer and the pH.

The interesting part is the organism clause, and it is interesting
because of what happens when we falsify it. We ran the plan a second
time with the organism constraint changed to an organism that does not
host MT-X. The \texttt{mt\_x} step went \vempty{} and every downstream
step went \vstarve{}. The plan returned nothing, and it returned
nothing \emph{with a blocker naming the corpus}.

\begin{principle}[a constraint is only a constraint if the answer
depends on it]
\label{prin:counterfactual}
For each constraint in a question, falsifying it must change the
emitted payload. A constraint that can be falsified while the payload
is unchanged is not being enforced; it is a comment on the plan.
\end{principle}

We applied \Cref{prin:counterfactual} to every constraint in every
plan. It caught a real defect in our own work, in \Cref{sec:g7}.

\subsection{Q3: substrate scope and product range}
\label{sec:q3}

\begin{quote}
\emph{What is the substrate scope and product range of the
Baeyer--Villiger monooxygenase BVMO-Y?}
\end{quote}

Three requests, three steps, all \vans: two reactions, two substrates,
two products. Every step succeeded and the rows are correct.

And the plan does not answer the question. This is \Cref{ex:scope}. The
corpus records substrates BVMO-Y \emph{was run on}; the question asks
for the substrates it \emph{would accept}. The plan therefore emits

\begin{lstlisting}[style=hfq]
emit substrates with provenance
     as extension of "substrate scope" because induction
\end{lstlisting}

and the emitted document carries, beside the verdicts:

\begin{lstlisting}[style=hfq]
"admissibility": {
  "asked_for": "substrate scope",
  "returned": "recorded extension",
  "answers_question": false,
  "reason": "... the inductive step from observed instances to
     scope is chemistry, not inference, and asserting it in an
     ontology would launder an assumption into a derived fact." }
\end{lstlisting}

Note that \texttt{verdict} remains \vans{} on every step while
\texttt{answers\_question} is \texttt{false}. Both are true. The
retrieval succeeded; the question was not answered. Collapsing these
into a single field would force us to lie about one of them.

\subsection{Q4: expected products under stated conditions}
\label{sec:q4}

\begin{quote}
\emph{What are the expected products of a biocatalytic kinetic
resolution reaction with enzyme PFE at pH 9 in HEPES buffer?}
\end{quote}

We flagged this in \Cref{sec:intro} as the one for a chemist to read
first. Here is the execution, verbatim from the evidence table:

\begin{table}[H]\centering
\begin{tabular}{llrl}
\toprule
step & verdict & $n$ & blocker \\
\midrule
\texttt{pfe\_reactions} & \vans     & 1 & --- \\
\texttt{pfe\_runs}      & \vans     & 1 & --- \\
\texttt{conditions}     & \vans     & 1 & --- \\
\texttt{at\_ph9}        & \vempty   & 0 & \emph{none} \\
\texttt{in\_hepes}      & \vstarve  & 0 & \bcorp \\
\bottomrule
\end{tabular}
\caption{Q4. Three requests issued. The last two rows are the result.}
\label{tab:q4}
\end{table}

Read the last two rows carefully, because they are two different
things and the difference is the paper.

\texttt{at\_ph9} is \vempty{} \textbf{with no blocker}. The corpus
holds a PFE kinetic resolution and holds its conditions, and those
conditions are not pH 9. The apparatus worked perfectly; the answer is
that the experiment at this pH is not in the record. That is a fact
about the chemistry, and \Cref{prin:noblocker} is why the field is
absent rather than null.

\texttt{in\_hepes} is \vstarve{} with blocker \bcorp. It was not
attempted. Its input was empty, so it was unreachable. This is a fact
about the plan's structure, not about HEPES: we learn nothing about
whether the corpus records HEPES buffers, because we never asked.

Now consider what the received pipeline returns for this question. An
empty solution sequence. One bit. The chemist cannot tell whether PFE
was never run under kinetic resolution, whether it was run but not at
pH 9, whether it was run at pH 9 but not in HEPES, or whether the
buffer vocabulary simply is not recorded --- and those four have four
different next actions, one of which is \emph{run the experiment}.

\begin{proposition}[a reasoner makes Q4 worse]
\label{prop:q4worse}
Suppose the ontology carries axioms sufficient to predict products of a
kinetic resolution from substrate and enzyme class. Then a reasoner
materialises product triples for the PFE reaction, and the query
returns them. The returned products are predictions under an assumed
mechanism, in the same vocabulary and the same graph as recorded
products, and the result set does not distinguish them.
\end{proposition}

This is \Cref{prop:launder} instantiated on the question the chemist
most wants answered, and it is where the two methods diverge most
sharply. The reasoner's answer is a list of products. Ours is: the
corpus has the reaction, has the conditions, and the conditions are not
the ones you named --- reported as \vempty{} without a blocker, which
is to say \emph{this is your chemistry, not our plumbing}. The
reasoner's list is more useful right up until it is wrong, and it
carries nothing that lets the reader find out which.

\subsection{Q5: instrument and wavelength for a dated observation}
\label{sec:q5}

\begin{quote}
\emph{With which device was the UV spectrum monitored during
biocatalytic transformation BT3 on March 23rd by Yuliia Dikova, and
which wavelength was monitored?}
\end{quote}

Two requests, four steps, all \vans: two monitored activities, of which
one is BT3, one on the date, one by the operator. The plan returns the
device and the wavelength.

Three constraints here are non-trivial --- transformation identity,
date, operator --- and we falsified each in turn. Changing the date
sends \texttt{on\_date} to \vempty{} and \texttt{by\_ydikova} to
\vstarve{}; changing the operator sends \texttt{by\_ydikova} to
\vempty{}. Each constraint is on the path from the question to the
payload.

\subsection{An implementation defect the framework caught in itself}
\label{sec:selfcatch}

While building the sources for Q5 we found that a predicate expressing
\emph{monitored on date $d$} matched nothing, for every date. The cause
was in our own adapter: the graph pattern applied its literal
restriction to the \emph{subject} position, so a literal appearing as
an object matched nothing at all. The step reported \vempty{} --- and
\vempty{} carries no blocker, which is to say the framework told us
\emph{the corpus has no such thing}, confidently, and was wrong.

We report this because it is the framework's own central failure mode
reproduced inside its implementation, and because it bounds what
\Cref{prin:noblocker} can promise. A blockerless \vempty{} means
\emph{the request completed and returned nothing}. It does not mean
\emph{the request asked what you meant}. That second guarantee is not
available at this layer, from us or from anyone, and a reader who takes
\vempty{} as evidence of absence is taking on the correctness of the
adapter's lowering as an assumption. The remedy we actually used was
\Cref{prin:counterfactual}: the defect surfaced because a constraint we
expected to matter did not change the payload when falsified.

% =====================================================================
\section{Eight dataset-discovery patterns}
\label{sec:generic}

The second group of questions is generic Chem-DCAT-AP dataset
discovery: \emph{give me all datasets $D$ about\ldots}. These are the
patterns a catalogue is supposed to serve, and they are a fairer test
of the standard than the bespoke questions, because they are what it
was designed for.

\begin{table}[H]\centering\small
\begin{tabular}{llrrl}
\toprule
 & pattern & req. & rows & outcome \\
\midrule
G1 & about a substance with compound C            & 2 & 2 & \vans \\
G2 & generated by activity of type T evaluating C & 2 & 2 & \vans{} + \textsf{vocabulary} \\
G3 & \ldots with Bruker spectrometer set to X     & 0 & --- & refused statically \\
G4 & about reaction R with product P              & 4 & 2 & \vans \\
G5 & \ldots with a starting material of compound C& 4 & 2 & \vans \\
G6 & \ldots with a solvent of compound C          & 3 & 2 & \vans{} + \textsf{vocabulary} \\
G7 & \ldots with a catalyst of compound C         & 5 & 2 & \vans \\
G8 & measured with a UV-vis spectrometer\ldots    & 0 & --- & refused statically \\
\bottomrule
\end{tabular}
\caption{The eight generic patterns. ``req.'' is
\texttt{requests\_issued}.}
\label{tab:generic}
\end{table}

\subsection{G3 and G8: refused at zero requests}
\label{sec:g8}

Both refusals have the same shape and G8 is the sharper one, so we take
it. The question asks for datasets measured with a UV-vis
spectrometer. The instrument registry \textsf{INST} declares
$\{\textsf{lookup}, \textsf{link}, \textsf{bind}, \textsf{batch}\}$: it
returns a record given a device identifier, and it accepts a batch of
them. The step requires $\{\textsf{pattern}, \textsf{neg}\}$, because
\emph{which of your devices are of kind K} is a pattern match over the
registry's contents, not a lookup by key. The check fails, and
\Cref{lst:refusal} is emitted with \texttt{requests\_issued} $= 0$.

Here is the part that matters. \textbf{The data to answer this question
exists in the federation.} The corpus contains device
\texttt{DEV:UV1900i} and activity \texttt{ACT:BT3}, and Q5
(\Cref{sec:q5}) successfully retrieves that very device. We are
refusing a question whose answer is present.

That is not a defect, it is the design working. Q5 asks \emph{which
device did this activity use}, which is a link from a known activity.
G8 asks \emph{which of your devices are UV-vis spectrometers}, which is
a classification over the registry's extension. The registry can do the
first and cannot do the second; the fact that one particular device
happens to be reachable by the first route does not make the second
route available. A pipeline that answered G8 by enumerating the devices
it happened to have seen would return a subset of unknown coverage,
formatted identically to a complete answer.

\begin{principle}[the refusal is about capability, not about data]
\label{prin:capnotdata}
A refusal states that no source can evaluate the constraint. It makes
no claim that the answer is absent. These are independent, and G8
exhibits a case where the first holds and the second does not.
\end{principle}

\Cref{prin:whole} is also visible here: G8's other steps are perfectly
satisfiable, and none of them ran. Had they run, the caller would hold
datasets satisfying the compound constraint and \emph{not} the
instrument constraint, in a payload whose shape does not distinguish
them from datasets satisfying both.

\subsection{G6: a vocabulary the corpus does not have}
\label{sec:g6}

G6 asks for datasets about a reaction whose \emph{solvent} has compound
C. Every step answers; two datasets are returned; and the plan carries
a \textsf{vocabulary} admissibility statement, because the corpus does
not record a solvent relation at all. It records reaction media. The
returned rows are about media, and media are not solvents.

The tempting repair is an ontology axiom: declare buffer a subproperty
of solvent, and the reasoner supplies the missing relation. This is
worth taking seriously because it is what the received method would
actually do, and it is wrong on the chemistry. The solvent of a
\SI{50}{\milli\Molar} Tris buffer is water; the buffer is a solute
system dissolved in it. An axiom equating the two produces triples
asserting that the solvent of the reaction is Tris, and those triples
enter the graph indistinguishable from observations
(\Cref{prop:launder}). The catalogue then answers G6 --- fluently, and
falsely.

\Cref{tab:generic} marks G2 with the same \textsf{vocabulary} gap, and
the two are not actually the same case: G2's subject is present in the
store and lacks a type, whereas G6's relation is absent entirely.
Our gap taxonomy has three values and does not separate these. We
record that as a limitation of the present partition rather than
claiming the three-way split is complete.

\subsection{G7, and a defect in our own plan}
\label{sec:g7}

G7 asks for datasets about a reaction whose \emph{catalyst} has
compound C. Five steps, five requests, all \vans, two datasets. The
answer is correct.

Our first version of this plan was also all-\vans{} and also returned
those two datasets, and it was wrong. It had a step that identified the
catalyst and checked its compound, and then a final step that retrieved
datasets bound to the \emph{reactions found before the check} rather
than to the checked set. Every verdict was \vans. The rows were right.

We found it by \Cref{prin:counterfactual}. Falsifying the catalyst
identity --- substituting an enzyme that is not the catalyst of that
reaction --- sent the \texttt{named} step to \vempty{}, exactly as it
should, and the final step \emph{still emitted two datasets}. The plan
was answering \emph{datasets about this reaction}, dressed as an answer
to \emph{datasets about this reaction with this catalyst}.

The repair is structural, and the plan carries the reasoning in a
comment because the next person to edit it needs it:

\begin{lstlisting}[style=hfq]
  # The reactions that survive the check, reached FROM the confirmed
  # catalyst. Binding `cat_rxns` here instead would make the check
  # decorative: the counterfactual run beside this one falsifies the
  # catalyst, `named` goes empty, and a plan that then binds
  # `cat_rxns` still emits two datasets -- a correct answer to
  # "datasets about this reaction" dressed as an answer to "datasets
  # about this reaction with THIS catalyst". Every constraint must sit
  # on the path from the question to the payload, or it is not a
  # constraint, it is a comment.
  let confirmed = from RXN
      ask catalysed_reactions(?c)
      with ?c in named
      within 20
\end{lstlisting}

\begin{principle}[a checked constraint that nothing depends on is
invisible]
\label{prin:invisible}
Verdicts classify what happened to each request. They cannot detect
that a step's output is unused downstream, because that is a property
of the plan's dataflow, not of any request. Only falsification exposes
it.
\end{principle}

We report this at length because it is a limit on our own claims. The
verdict algebra is not a correctness check on plans. It tells you what
happened to each step; it does not tell you that the steps compose into
the question you asked. The counterfactual discipline of
\Cref{prin:counterfactual} is the part that does that work, and it is
methodology, not machinery --- it must be performed, and nothing in the
executor enforces it.

\subsection{One further contrast worth stating}

G6 and G7 are the same question shape --- \emph{datasets about a
reaction one of whose participants has compound C} --- differing only
in the participant role. G7 answers cleanly and G6 does not, and the
entire difference is that someone recorded a catalyst relation and
nobody recorded a solvent relation.

That is worth saying plainly to a catalogue designer. The difference
between a question this infrastructure answers and one it cannot is
not depth of reasoning, expressiveness of the ontology, or power of the
engine. It is whether the relation was written down. Every layer above
the record inherits that, and a layer that appears to escape it ---
by deriving the missing relation from an axiom --- has not escaped it,
it has stopped reporting it.

% =====================================================================
\section{Two design decisions that look like omissions}
\label{sec:omissions}

Two things a reader coming from SPARQL will expect to find, and will
not. Both are deliberate, and both follow from \Cref{obs:onebit} rather
than from expedience.

\subsection{There is no set difference in the plan language}
\label{sec:nominus}

The plan language has \texttt{union}, \texttt{intersect}, \texttt{join
\ldots on}, \texttt{filter \ldots where}, and \texttt{map \ldots via}.
It has no \texttt{minus}. SPARQL has \texttt{MINUS} and \texttt{FILTER
NOT EXISTS}, and their absence here is the single largest expressive
difference between the two languages.

The reason is that a plan-level difference operator would let a plan
manufacture a negative answer out of two positive ones, and thereby
destroy exactly the distinction the verdicts exist to preserve.

\begin{proposition}[a plan-level difference erases starvation]
\label{prop:nominus}
Let $A$ and $B$ be steps, and suppose a plan may write $A \setminus B$
as a plan-level operation over their payloads. If $B$ is \vstarve{} ---
unreachable, never issued --- then its payload is empty and $A
\setminus B = A$. The combined step is therefore indistinguishable, in
both payload and cardinality, from the case where $B$ was reached and
genuinely matched nothing.
\end{proposition}

\begin{proof}
Immediate: $A \setminus \emptyset = A$ regardless of why $B$'s payload
is empty, and the plan-level operator sees only payloads. The verdict on
$B$ survives in the step record, but the \emph{payload} the caller reads
carries no trace of it, and any downstream step binding the difference
inherits the erasure rather than the verdict.
\end{proof}

This is \Cref{obs:onebit} reappearing one level up, inside our own
language, which is why we regard it as disqualifying rather than merely
inconvenient. The alternative we adopt is that negation is a
\emph{source capability}, \textsf{neg}, declared by a source that can
actually evaluate an absence over its own extension. A step requiring
\textsf{neg} either lands on a source that has it --- in which case the
absence is evaluated where the data live, by something that knows the
extent of what it is quantifying over --- or the plan is refused. G8 is
this case: the instrument registry does not declare \textsf{neg}, and
the plan is refused at zero requests (\Cref{sec:g8}) rather than
answered by differencing two things the registry \emph{can} return.

The cost is real and we state it. Questions that a SPARQL author would
write with \texttt{FILTER NOT EXISTS} over a single materialised graph
must here be either pushed into a source that declares \textsf{neg}, or
refused. On a federation where every source is a full SPARQL endpoint
over a complete graph, this buys nothing and costs expressiveness. On
the federation the chemist actually has --- REST sequence services,
keyed instrument registries, a local provenance store --- it is the
difference between a negative answer that means something and one that
means \emph{we did not look}.

\subsection{Comparability is a property of a pair, not of a dataset}
\label{sec:pairs}

Q2 and Q5 look like retrieval questions and are usually treated as such:
\emph{which buffer}, \emph{which device}. But the reason a chemist asks
them is almost never to learn a buffer identity in isolation. It is to
decide whether two runs may be compared --- whether this transformation
and that one were carried out under conditions close enough that a
difference in outcome is about the chemistry.

That target is not a property of either dataset. It is a relation on the
pair, and it is not one an ontology can store, for a reason that is
arithmetic rather than architectural: with $n$ runs there are
$\binom{n}{2}$ pairs, and the relation depends on a tolerance the asker
supplies at question time. A catalogue that pre-computed comparability
would have to fix the tolerance before knowing the question, and would
be storing a quantity that grows quadratically and is invalidated by
every new run.

We do not solve this here, and we flag it as the most promising place
for the plan layer to go next, because the plan layer is where it
naturally sits. A plan already carries steps whose inputs are bound by
earlier steps; a comparability step is one whose input is bound by
\emph{two}, with a tolerance as a parameter, and whose verdict would
distinguish \emph{these runs are not comparable} from \emph{the record
does not say whether they are comparable} --- which is again
\Cref{obs:onebit}, in the place where a chemist would most be misled by
its collapse. What blocks us is not the plan machinery but the
representation of the tolerance itself, which is the open problem
recorded under \Cref{sec:related}'s measurement paragraph.

% =====================================================================
\section{Full evidence}
\label{sec:evidence}

Every number quoted in this paper comes from the table below, which is
the executor's own per-step output over the thirteen plans. We give it
in full rather than in summary so that a reader can check the claims
against it. ``req.'' is \texttt{requests\_issued}; ``wc'' is
well-capability.

\begin{table}[H]\centering\small
\begin{tabular}{llrl}
\toprule
plan & step & $n$ & verdict \\
\midrule
\multicolumn{4}{l}{\textbf{Q1} \quad req.\ 3, halted false, wc true} \\
& \texttt{enzymes} & 5 & \vans \\
& \texttt{bacterial} & 4 & \vans \\
& \texttt{candidates} & 4 & \vans \\
& \texttt{keyed} & 3 & \vans{} \quad ($r_\mu = 0.75$, $a_\mu = 1.0$) \\
& \texttt{no\_cys} & 2 & \vans \\
\midrule
\multicolumn{4}{l}{\textbf{Q2} \quad req.\ 5, halted false, wc true} \\
& \texttt{bs\_enzymes} & 2 & \vans \\
& \texttt{mt\_x} & 1 & \vans \\
& \texttt{mt\_reactions} & 1 & \vans \\
& \texttt{runs} & 2 & \vans \\
& \texttt{settings} & 2 & \vans \\
\midrule
\multicolumn{4}{l}{\textbf{Q3} \quad req.\ 3, halted false, wc true} \\
& \texttt{reactions} & 2 & \vans \\
& \texttt{substrates} & 2 & \vans \\
& \texttt{products} & 2 & \vans \\
\midrule
\multicolumn{4}{l}{\textbf{Q4} \quad req.\ 3, halted false, wc true} \\
& \texttt{pfe\_reactions} & 1 & \vans \\
& \texttt{pfe\_runs} & 1 & \vans \\
& \texttt{conditions} & 1 & \vans \\
& \texttt{at\_ph9} & 0 & \vempty{} \quad (no blocker) \\
& \texttt{in\_hepes} & 0 & \vstarve{} \quad (\bcorp) \\
\midrule
\multicolumn{4}{l}{\textbf{Q5} \quad req.\ 2, halted false, wc true} \\
& \texttt{measured} & 2 & \vans \\
& \texttt{bt3} & 1 & \vans \\
& \texttt{on\_date} & 1 & \vans \\
& \texttt{by\_ydikova} & 1 & \vans \\
\midrule
\multicolumn{4}{l}{\textbf{G1} req.\ 2 \quad \texttt{rxns} 1, \texttt{datasets} 2 --- all \vans} \\
\multicolumn{4}{l}{\textbf{G2} req.\ 2 \quad \texttt{rxns} 1, \texttt{datasets} 2 --- all \vans{}, gap \textsf{vocabulary}} \\
\multicolumn{4}{l}{\textbf{G3} req.\ \textbf{0}, halted true, wc \textbf{false} \quad \texttt{brukers} \vsurf} \\
\multicolumn{4}{l}{\textbf{G4} req.\ 4 \quad \texttt{products}/\texttt{has\_p}/\texttt{rxns} 1, \texttt{datasets} 2 --- all \vans} \\
\multicolumn{4}{l}{\textbf{G5} req.\ 4 \quad \texttt{inputs} 2, \texttt{has\_c} 1, \texttt{rxns} 1, \texttt{datasets} 2 --- all \vans} \\
\multicolumn{4}{l}{\textbf{G6} req.\ 3 \quad \texttt{rxns} 1, \texttt{datasets} 2, \texttt{media} 2 --- all \vans{}, gap \textsf{vocabulary}} \\
\multicolumn{4}{l}{\textbf{G7} req.\ 5 \quad \texttt{cat\_rxns}/\texttt{catalysts}/\texttt{named}/\texttt{confirmed} 1, \texttt{datasets} 2 --- all \vans} \\
\multicolumn{4}{l}{\textbf{G8} req.\ \textbf{0}, halted true, wc \textbf{false} \quad \texttt{uvvis} \vsurf} \\
\bottomrule
\end{tabular}
\caption{Per-step evidence for all thirteen plans.}
\label{tab:evidence}
\end{table}

Three facts in \Cref{tab:evidence} deserve to be read off explicitly,
because they are the paper's empirical content and a summary would hide
them.

First, \texttt{Q1.keyed} is the \emph{only} row in the entire suite
carrying retention and amplification. Where no namespace is crossed
there is no attrition to measure and the fields are absent. We say this
rather than leaving a reader to infer that map metrics are pervasive,
which they are not.

Second, G3 and G8 are the only rows with \texttt{requests\_issued} $=
0$, and they are exactly the rows with \texttt{well\_capability} false.
That correspondence is \Cref{cor:refuse} holding operationally: the
static check is the thing that halts, and it halts before contact.

Third --- and this is the one worth dwelling on --- \textbf{eleven of
the thirteen plans have every step \vans}. The verdict algebra spends
almost all of its time saying \emph{nothing went wrong}. That is not a
weak result, it is the correct shape for a working federation, and it is
worth stating because a reader could otherwise take this paper as
claiming that federated retrieval is mostly broken. It is not. The
apparatus earns its keep on the two rows in Q4 and the two refusals, and
those five rows out of the roughly forty in the table are precisely the
ones a one-bit interface would have rendered as a single undifferentiated
absence.

% =====================================================================
\section{Validation, and what it does not establish}
\label{sec:validation}

Sixteen checks (V1--V16) run over the compiler and the executor. They
cover: that the verdict assignment is total and deterministic; that the
blocker map is partial exactly at \vempty; that a plan failing the
capability check issues no request; that allocation is feasible and
exhausts the budget when demand exceeds it; that no step is dropped by
rule (\Cref{cor:nodrop}); that retention and amplification agree with
\Cref{def:retamp} on constructed maps; that the blame chain terminates;
that emitted documents round-trip through JSON with the blocker field
\emph{absent} rather than null; and that the parser rejects the
constructs it is specified to reject. All sixteen hold.

The round-trip check deserves a note, because it is the least
interesting-looking and guards the most. \Cref{prin:noblocker} says the
blocker map is partial at \vempty. A serialiser that emitted
\texttt{"blocker": null} would satisfy every other check while quietly
reintroducing a distinguished absent value --- the exact move that
\Cref{sec:related} criticises SQL's NULL for. The check asserts the key
is missing, not that its value is null.

\begin{remark}[the disclaimer that matters more than the result]
\label{rem:validation}
V1--V16 test the implementation against the theory. They do not test the
theory against the world. That a verdict is assigned deterministically
says nothing about whether it is the right verdict; that no request is
issued for an ill-capable plan says nothing about whether the capability
declaration was true. The check that would matter most --- does a
source's declaration match its behaviour --- is precisely the one we
cannot run, for the reason in \Cref{sec:noprobing}.
\end{remark}

There is a second gap the suite cannot reach, and \Cref{sec:g7} is the
evidence for it: no check detects a plan whose steps all succeed and
which answers a different question. That defect is invisible to
everything in V1--V16, and it was found by hand. We would rather report
sixteen checks with both caveats attached than a larger number that
implies coverage we do not have.

% =====================================================================
\section{On not probing a third party's endpoint}
\label{sec:noprobing}

The evaluation runs against a local fixture, constructed to mirror the
structure of the group's records --- the same sources, the same
namespace split, the same sparse recording --- rather than against a
live public endpoint. Two reasons, and the second is the binding one.

The practical reason is reproducibility. A paper whose numbers depend on
a third party's deployment state reports something unrepeatable: the
endpoint changes and \Cref{tab:evidence} becomes unverifiable.

The reason we regard as binding is ethical. Characterising a source's
\emph{true} capability set --- which is what \Cref{rem:honesty} says we
cannot verify --- would require issuing requests designed to elicit
internal behaviour: probing for materialisation caps, for silent
truncation, for the point at which an aggregate stops being the
aggregate. That is adversarial traffic against a public research service
maintained on a shared budget, and it is not ours to send. We hold to
this even though it leaves the largest gap in the design unclosed. Given
the choice between naming a gap and closing it by that route, we name it.

There is a related observation we can report but not explain. On
10 August 2026, two requests to a public endpoint, byte-identical except
for one line, returned 2 and 397 rows respectively --- a ratio of
$198.5$. We do not know why. Every hypothesis we can form is a
hypothesis about the deployment's internals, and discriminating among
them requires exactly the probing above. So it stands here as an
undiagnosed observation rather than a characterised behaviour. We think
reporting it undiagnosed is more useful than either omitting it or
speculating, and we note that it is the kind of thing
\Cref{prop:under} is a response to: a client that cannot rule out
truncation should decline to declare the capability, not guess at the
threshold.

The implementation reflects the constraint in code rather than in
intention. The adapter module carries it in its docstring: every adapter
resolves against a local fixture or a local engine; no adapter performs
network I/O, and none may be added that does.

% =====================================================================
\section{Related work}
\label{sec:related}

\paragraph{Source capability description.} Describing what a source can
do is old and well studied: mediator architectures
\citep{Wiederhold1992}, source descriptions for query answering
\citep{Levy1996}, and capability-aware optimisation across heterogeneous
sources \citep{Haas1997}. SPARQL federation inherits the problem
\citep{Prudhommeaux2013}, systems such as FedX \citep{Schwarte2011} plan
around source capability, and the space is surveyed by
\citet{Saleem2016}. Nothing in \Cref{sec:calculus} is novel as
\emph{description}, and we do not claim it is.

The difference is what happens to a capability mismatch. In all of the
above it is a planner-internal fact: the planner routes around it,
rewrites, or fails, and the caller receives a result set. Here it is a
first-class output --- named, with the required set, the declared set,
and the difference, at \texttt{requests\_issued} $= 0$. We are not
claiming a better planner. We are claiming that the mismatch is
information the caller needs, and that the conventional interface has
nowhere to put it, which is \Cref{obs:onebit}.

\citet{Aranda2013} is close in spirit: it surveys the reliability of
public SPARQL infrastructure and finds it uneven. \Cref{prop:under} can
be read as asking what a client should do about that finding, and
answering: declare the capability away, and accept refusals in exchange
for never receiving a truncated aggregate that has the shape of a
complete one.

\paragraph{Nulls and certain answers.} The relational literature has
long distinguished the reasons a value or a row may be absent
\citep{Codd1979,Imielinski1984}, and \citet{Libkin2016} shows how far
SQL's three-valued treatment sits from certain-answer semantics. Our
\Cref{def:rules} is a coarser instrument aimed at a different level: not
\emph{is this value unknown} but \emph{did this step get to look}. The
connection is that both are refusals to let one distinguished value ---
NULL there, the empty solution sequence here --- carry more meanings
than it can separate. The debt is direct, and \Cref{prin:noblocker} is
the place we try hardest to honour it: rather than introducing a
distinguished ``no blocker'' value, which would recreate the problem, we
make the field absent.

\paragraph{Reasoning.} We use description-logic reasoning
\citep{Baader2003,Grau2008,Motik2009} as the foil throughout, and we
want to be precise about the criticism, because it is narrower than the
framing might suggest. Nothing here says a reasoner computes the wrong
entailments; under its axioms it computes exactly the right ones. The
criticism is about \emph{materialisation into a shared vocabulary}:
entailed and asserted triples become indistinguishable to any query over
the result (\Cref{prop:launder}), and on sparse experimental data the
entailed fraction can be large while its axioms go unexamined by the
person reading the answer. The closed-world and default-reasoning
literature \citep{Reiter1978,Reiter1980} treats exactly this as
something to mark rather than absorb, and our admissibility statement is
a coarse, plan-level version of that marking --- coarser than a default
theory, and correspondingly cheaper to write and to read.

\paragraph{The standard and its domain.} Chem-DCAT-AP extends DCAT-AP
\citep{DCATAP2015,DCAT2020} with catalysis vocabulary, in the FAIR
tradition \citep{Wilkinson2016} and the biocatalysis standardisation
effort around it \citep{Doerr2020,Pleiss2021,Bornscheuer2012}; the
sources it federates are the usual ones
\citep{UniProt2023,Wittig2018,Chang2021,Hastings2016}, with provenance
in PROV-O \citep{PROVO2013}, and a query layer specified by
\citet{Harris2013} over engines such as \citet{Pyoxigraph}. We repeat
that we keep all of this. The present paper changes the answering layer,
not the vocabulary.

\paragraph{Measurement.} One connection we flag without developing. The
metrology guide \citep{GUM2008} requires a measured value to carry its
uncertainty. RDF does not: a triple asserting \texttt{pH 9.0} and a
triple asserting \texttt{pH 9.0} $\pm$ \texttt{0.5} are structurally the
same triple with different literals, and no query over them can be made
to respect the difference. Q4 sits directly on this --- \emph{at pH 9}
is a sharp cut through a quantity that was measured with a tolerance ---
and so does the comparability question of \Cref{sec:pairs}, which is
\emph{entirely} a question about tolerances. We do not solve it here. We
note that the framework this work builds on handles it by making every
value carry a strictly positive floor, so that a sharp cut is
\emph{untypeable} rather than merely discouraged
\citep{HegelFederated2026}, and that reconciling a typed floor with RDF
literals is open.

% =====================================================================
\section{Limitations}
\label{sec:limits}

Beyond the four limits stated in \Cref{sec:intro}, three more have
emerged from doing the work, and they are sharper than the ones we knew
to state in advance.

\begin{enumerate}[label=(\roman*),leftmargin=2.2em]
\item \textbf{The verdicts do not check plans.} \Cref{sec:g7} is our own
  defect: every step \vans, the right rows, the wrong question. The
  algebra classifies what happened to each request, and a constraint
  whose output nothing consumes is a dataflow property that no verdict
  can see. The remedy we used --- counterfactual falsification of every
  constraint (\Cref{prin:counterfactual}) --- is methodology. Nothing in
  the executor enforces it, and a plan author who skips it gets no
  warning. A static check that every non-terminal step's output is
  transitively consumed by the emitted step would catch this particular
  shape, and we have not built it.
\item \textbf{The gap taxonomy is under-resolved.} Three values, and
  \Cref{sec:g6} shows \textsf{vocabulary} covering two structurally
  different situations --- a missing type on a subject that is present,
  and a relation that is absent from the schema entirely. A finer
  partition would be more informative and harder to assign consistently.
  We have not found the right cut, and we would rather ship three
  honest values than five that the plan author assigns by coin flip.
\item \textbf{A blockerless \vempty{} inherits the adapter's
  correctness.} \Cref{sec:selfcatch} is the demonstration: our own
  lowering asked the wrong question, and the framework reported
  \emph{nothing there}, confidently. \vempty{} means the request
  completed and returned nothing. It does not mean the request expressed
  the constraint the plan author intended. That second guarantee is not
  available at this layer, and a reader who treats \vempty{} as evidence
  of absence is taking the adapter's lowering on trust.
\end{enumerate}

\section{Conclusion}
\label{sec:conclusion}

Thirteen questions from working biocatalysis practice, put to a plan
language rather than to a query language. Eight return rows. Two are
refused before any request is issued, one of them for data that
demonstrably exists in the federation --- and that refusal is correct,
because no source can evaluate the constraint and a partial answer would
have had the shape of a complete one. Two return rows under a named
admissibility gap. One returns \vempty{} without a blocker and a
downstream \vstarve{} with one, which is the distinction the whole
apparatus exists to preserve.

We keep the ontology. The change is at the answering layer, and the
argument for it is a single observation: a result set carries one bit of
failure information, and laboratory catalysis data --- partial by
nature, federated by necessity, unevenly recorded by circumstance ---
routinely produces empty results for reasons that one bit cannot
separate. Distinguishing them is not a matter of writing better queries,
because the distinction is a property of the execution and the interface
has no field for it.

What we can promise is narrower than the tally suggests, and we would
rather end on that than on the tally. The apparatus reports faithfully
what happened to each request. It does not certify that the requests
expressed the question --- our own G7 plan returned the right rows for
the wrong reason, and our own adapter reported an empty corpus that was
not empty. It does not verify a source's capability declaration, and it
deliberately declines the only route that could. What it does is smaller
and, we think, worth having: when this infrastructure fails to answer a
chemist's question, the document it returns says whether the failure was
in the chemistry or in the plumbing. That is one bit more than the
alternative, and it is the bit the chemist needs in order to know
whether the next move is to fix a cross-reference or to run the
experiment.

\bibliographystyle{plainnat}
\bibliography{references}

\end{document}
